\section{The persistence problem}

Living organization persists while its material components are replaced.
Enzymes are degraded and resynthesized, nucleosomes are diluted at every
replication fork, and cells in a tissue turn over while the tissue continues to
perform its role. The motivating question is how the consequence of an
interaction can persist after the object or arrangement that first carried it
has changed or disappeared, without slipping between incompatible meanings of
``information.'' None of these examples requires information to float free of
physical realization. They distinguish persistence beyond an \emph{original}
carrier from existence without \emph{any} carrier.

Embodied-information accounts in theoretical biology already treat patterns as
constraints on physical processes and emphasize their historical accumulation
\citep{Farnsworth2022Information}. Constraint-closure theory models conserved
structures that act on processes while depending on one another for their
maintenance \citep{MontevilMossio2015Closure}, and the theory of the adjacent
possible treats organization as enabling futures that were previously
inaccessible \citep{Kauffman2000Investigations,Kauffman2014Prolegomenon}. A new
proposal must therefore earn its place through operational criteria, not
through the language of embodiment or constraint alone.

We state the proposal as follows:
\begin{quote}
\textbf{Information-Constraint Principle (ICP).} Physically effective
information is interaction-formed, embodied constraint on an explicitly indexed
family of reachable or probable futures; it persists when that constraint role
is lawfully re-instantiated beyond its original carrier.
\end{quote}
ICP arose within a broader information-first modelling program whose
artificial-life testbed, IOTA, is reported separately
\citep{Clark2026IotaConstraint}. The present paper does not depend on that
program's physical ontology. It asks a narrower question: does the combination
of historical formation, physical embodiment, sufficiency for a declared
family of future responses, causal carrier manipulation, succession across
replacement of the original carrier, and robustness to representation yield a
testable claim that existing theories do not already make?

The phrase ``constrains the future'' is easy to satisfy. Given a complete
physical state and law, virtually every state difference can alter a future
trajectory. Control theory studies reachable sets; computational mechanics
conditions futures on pasts; and semantic-information work erases correlations
and measures effects on later viability. If ICP merely renames these objects it
may be useful language, but it is not a new principle. The paper is
accordingly adversarial: it formalizes ICP (Section~\ref{sec:formalism}), tests
it against its closest native formalisms (Section~\ref{sec:redundancy}),
applies it to a well-characterized biological memory
(Sections~\ref{sec:biology}--\ref{sec:ino1}), and grades what survives
(Sections~\ref{sec:audit}--\ref{sec:path}).
